\scanpage{236}
% morphism of complexes
On trouve alors un homomorphisme
\[
 \varphi:\mathcal L^\bullet\longrightarrow\mathscr H^\bullet
\]
par
\diagram{\begin{tikzcd}[ampersand replacement=\&,column sep=small,row sep=large]
 \cdots\arrow[r]\&0\arrow[r]\&\unsure{\mathscr G|\mathcal G}\arrow[r,"i"]\arrow[d,"{\deleted{\unsure{k}}\varphi^{-1}}"]\&
 \overline{\mathscr G}\arrow[r]\arrow[d,"{\varphi^0}"]\&0\arrow[r]\&\cdots\\
 \&\mathscr H^{-2}\arrow[r]\&\mathscr H^{-1}\arrow[r,"{\delta^{\unsure{-1}}}"]\&
 \mathscr H^0\arrow[r]\&\mathscr H^1\arrow[r]\&\cdots
\end{tikzcd}}
défini par les conditions
\[
 \varphi^{-1}=\unsure{k|h}\qquad \varphi^0|\unsure{\mathscr P|\mathcal P}=u
\]
et la commutativité du carré

% explanation of a homomorphism out of the extension
[\deleted{Du fait que} \inserted{Un} hom de $\overline{\mathscr G}$ dans un un faisceau $F$ $=$ \{ Hom $v$ de $\mathscr G^{\unsure{\cdot\cdot}}$ dans $F$, et hom $u$ de $\mathscr P$ dans $F$, avec condition $u(pg)=u(p)v(g)$\}

Donc prenant \deleted{$\delta\varphi^{-1}$} $F=\mathscr H^0$, et $u,v$ comme dessus, on trouve $\varphi^0$, avec $\varphi^0 i=u$. Mais $u=\delta\deleted{\unsure{\varphi^{-1}|\varphi}}k=\delta\varphi^{-1}$ ... ]

% tensoring, preliminary composite and boxed result
Tensorisant avec $\underline{\mathcal O}_X$, on trouve
\diagram{\scriptsize\begin{tikzcd}[ampersand replacement=\&,column sep=small]
 \underline{\mathcal O}_X\arrow[r,"{\text{dans cat.}}","{\text{dérivée}}"']\&
 \deleted{\unsure{\mathcal L^0|\mathscr L}}\inserted{\unsure{\mathcal L^\bullet|\mathscr L}}\arrow[r]\&\mathscr H\arrow[r]\&
 \mathbb R\underline{\mathrm{Hom}}(K',K'')
\end{tikzcd}}
\deleted{On trouve donc,} donc
\[
 \boxed{\underline{\mathcal O}_X\unsure{\xrightarrow{\sim}|\longrightarrow}
 \mathbb R\underline{\mathrm{Hom}}(K',K'')\quad\text{ds cat. dérivée}}
\]
Donc
\[
 \deleted{\mathbb R\Gamma_X(\unsure{\mathcal L})}\simeq
 \mathbb R\Gamma_X(\deleted{\unsure{\mathbb Z}}\underline{\mathcal O}_X)\longrightarrow
 \deleted{\mathbb R\Gamma_X\unsure{\mathscr H}}\longrightarrow\mathbb R\mathrm{Hom}(K',K'')
\]
d'où
\[
 \boxed{\mathbb R\Gamma_X(\deleted{\unsure{\mathbb Z}}\underline{\mathcal O}_X)\longrightarrow\mathbb R\mathrm{Hom}(K',K'')}
\]
Par suite, \deleted{faisant} passant au $\mathbb R\Gamma^0_X$, on trouve un \unsure{él\textsuperscript{mt}|élémt|élément} de $\mathbb R^0\mathrm{Hom}(K',K'')$, i.e. un hom de la cat\unsure{e|\textsuperscript{e}|é|.} dérivée
% archive sheet number: 235
